\documentclass[11pt]{article}
\usepackage[solutions]{../../../mystyle/Style}
\usepackage[margin=1in]{geometry}
\usepackage{graphicx}
\usepackage{booktabs,longtable,array,listings}
\usepackage[colorlinks=true,linkcolor=blue,urlcolor=blue]{hyperref}
\lstset{language=[LaTeX]TeX,basicstyle=\small\ttfamily,columns=fullflexible,
  keepspaces=true,breaklines=true,showstringspaces=false,frame=single,
  rulecolor=\color{black!20},backgroundcolor=\color{black!3},tabsize=2}
\setlength{\emergencystretch}{2em}
\title{Writing LaTeX\\\large A Practical Guide and Personal Snippet Reference}
\author{\defaultauthor}
\date{\today}
\begin{document}
\maketitle
\noindent Use this guide while writing: read the explanation, copy the source,
then compare it with the rendered example. The final sections document all
117 personal VS Code snippets and the commands in your current style file.
\tableofcontents
\newpage

\section{From a blank file to a document}
LaTeX turns a plain-text \texttt{.tex} source file into a typeset PDF.
Save your edits, compile, and read the PDF and any compiler diagnostics.
Changes in the source appear in the PDF after a successful compilation.
Start with this minimal document:
\begin{lstlisting}
\documentclass{article}
\title{My Notes}
\author{Hem Dadlani}
\date{\today}

\begin{document}
\maketitle
\section{Introduction}
My first paragraph. Inline math: $x^2+1$.
\end{document}
\end{lstlisting}
The \emph{preamble} is everything before \lstinline|\begin{document}|:
choose the document class, load packages, define commands, and set metadata
there. The body holds the visible content. Setting a title does not print it;
\lstinline|\maketitle| prints it.

\subsection{Using your style file}
Your style loads the usual math packages, theorem environments, homework
helpers, and shortcuts such as \lstinline|\R| and \lstinline|\abs{x}|.
For a file in this guide's folder, use:
\begin{lstlisting}
\documentclass{article}
\usepackage[solutions]{../../../mystyle/Style}
\usepackage[margin=1in]{geometry}
\usepackage{graphicx} % for images
\usepackage[hidelinks]{hyperref} % clickable references
\end{lstlisting}
The path depends on your file's location: each \texttt{../} goes up one
folder. The original reference document, one folder above this guide, uses
\texttt{../../mystyle/Style}. Templates currently insert
\texttt{../../../mystyle/Style}; check that path whenever starting somewhere
new. Filename capitalization matters on some systems.
Do not paste a complete document template into another document's body.

\section{Writing text and organizing pages}
A command normally begins with a backslash. Braces contain required arguments;
square brackets often contain optional arguments. An environment has matching
\lstinline|\begin{name}| and \lstinline|\end{name}| commands.
\begin{lstlisting}
\section{Groups}
\subsection{Examples}
\subsubsection{Finite examples}
This is \textbf{bold}, \emph{emphasized}, and \texttt{code}.

This starts a new paragraph. % This comment is not printed.
\end{lstlisting}
This is \textbf{bold}, \emph{emphasized}, and \texttt{code}.

A blank source line starts a paragraph. A single source line break normally
acts like a space. Use \lstinline|\\| for a deliberate line break in an
alignment or table; use blank lines for ordinary paragraphs.
\lstinline|\newpage| starts a page. \lstinline|\tableofcontents| creates the
contents from section headings; compile again if it is out of date.
\lstinline|\section*{Acknowledgments}| makes an unnumbered heading that does
not automatically enter the contents.

\subsection{Characters that need escaping}
These characters have special meanings in source:
\begin{lstlisting}
\%  \$  \&  \#  \_  \{  \}
\textbackslash{}  \textasciitilde{}  \textasciicircum{}
\end{lstlisting}
Result: \% \$ \& \# \_ \{ \} \textbackslash{} \textasciitilde{}
\textasciicircum{}. A percent sign starts a comment; an underscore denotes a
subscript in math. For literal commands in a document, use
\lstinline|\verb| with a delimiter, or a \texttt{verbatim} environment for a
block. These literal-code forms generally cannot go inside section titles.

\section{Math: choosing the right environment}
Use inline math inside a sentence, and display math for a separate centered
formula. Most mathematical commands must be inside one of these math modes.
\begin{lstlisting}
For $x\in\R$, we have $x^2\geq 0$.
\[
  \frac{a+b}{c}=\sqrt{x^2+1}
\]
\end{lstlisting}
For $x\in\R$, we have $x^2\geq 0$.
\[\frac{a+b}{c}=\sqrt{x^2+1}\]
Braces group an argument: \lstinline|x^{n+1}| raises the whole expression,
whereas \lstinline|x^n+1| only raises the \(n\).
\begin{lstlisting}
x_i,\quad x_{i+1},\quad x^{n+1},\quad
\frac{1}{n},\quad \sqrt{x},\quad \sqrt[3]{x}
\end{lstlisting}
\[x_i,\quad x_{i+1},\quad x^{n+1},\quad
\frac{1}{n},\quad \sqrt{x},\quad \sqrt[3]{x}\]

\subsection{Numbered equations and references}
\begin{lstlisting}
\begin{equation}
  \sum_{i=1}^n i=\frac{n(n+1)}{2}
  \label{eq:sum}
\end{equation}
See Equation~\eqref{eq:sum}.
\end{lstlisting}
\begin{equation}\sum_{i=1}^n i=\frac{n(n+1)}{2}\label{eq:sum}\end{equation}
See Equation~\eqref{eq:sum}. Your style numbers equations within sections.
An asterisk removes numbering: \texttt{equation*} is unnumbered. Your
\texttt{eq} snippet inserts \texttt{equation*}; remove the star for a number.
Do not nest \texttt{equation} or \texttt{align} inside dollar signs or
\lstinline|\[ ... \]|.

\subsection{Calculations on several lines}
\begin{lstlisting}
\begin{align*}
  (x+1)^2 &= x^2+2x+1 \\
          &= x(x+2)+1
\end{align*}
\end{lstlisting}
\begin{align*}(x+1)^2 &= x^2+2x+1 \\ &= x(x+2)+1\end{align*}
The ampersand marks the alignment column, and two backslashes end a row.
Use \texttt{align} for numbered rows, \texttt{align*} for unnumbered rows,
and \lstinline|\notag| to suppress an individual row's number.
\texttt{aligned} is a smaller alignment used \emph{inside} math mode.

\subsection{Words, spacing, and common symbols}
\begin{lstlisting}
f(x)=x^2 \quad \text{for all } x\in\R
A\subseteq B,\quad x\neq y,\quad x\leq y,\quad x\geq y
P\implies Q,\quad P\iff Q,\quad f:A\to B
\alpha,\beta,\gamma,\delta,\eps,\theta,\lambda,\pi,\sigma,\omega
\end{lstlisting}
\[f(x)=x^2 \quad \text{for all }x\in\R\]
\[A\subseteq B,\quad x\neq y,\quad x\leq y,\quad x\geq y\]
\[P\implies Q,\quad P\iff Q,\quad f:A\to B\]
Use \lstinline|\text{...}| for words in math. Ordinary typed spaces in math
are mostly ignored; \lstinline|\,|, \lstinline|\quad|, and
\lstinline|\qquad| insert increasing amounts of space.

\subsection{Sums, products, limits, integrals, and derivatives}
\begin{lstlisting}
\sum_{i=1}^n i,\quad \prod_{i=1}^n i=n!,\quad
\lim_{n\to\infty}\frac1n=0,\quad
\int_0^1 x^2\,dx=\frac13,\quad
f'(x),\quad \frac{df}{dx},\quad \frac{\partial f}{\partial x}
\end{lstlisting}
\[\sum_{i=1}^n i,\quad \prod_{i=1}^n i=n!,\quad
\lim_{n\to\infty}\frac1n=0,\quad \int_0^1x^2\,dx=\frac13\]
\[f'(x),\quad \frac{df}{dx},\quad \frac{\partial f}{\partial x}\]
Use operator commands such as \lstinline|\sin x| and \lstinline|\log x|
instead of typing their names as ordinary math letters.

\subsection{Delimiters, matrices, and cases}
\begin{lstlisting}
\left(\frac{a}{b}\right),\quad \abs*{\frac{x}{y}}
\begin{pmatrix} a & b \\ c & d \end{pmatrix}
f(x)=\begin{cases}
  x^2 & \text{if }x\geq 0 \\
  -x & \text{if }x<0
\end{cases}
\end{lstlisting}
\[\left(\frac{a}{b}\right),\quad \abs*{\frac{x}{y}},\quad
\begin{pmatrix}a&b\\c&d\end{pmatrix}\]
\[f(x)=\begin{cases}x^2&\text{if }x\geq0\\-x&\text{if }x<0\end{cases}\]
Put matrix and case snippets inside math mode. Use \texttt{bmatrix} for
square brackets. Match every \lstinline|\left| with a \lstinline|\right|;
\lstinline|\right.| supplies an invisible matching delimiter.

\section{Theorems, definitions, examples, and proofs}
\begin{lstlisting}
\begin{definition}[Even integer]
  An integer $n$ is even if $n=2k$ for some $k\in\Z$.
\end{definition}
\begin{theorem}\label{thm:even}
  The sum of two even integers is even.
\end{theorem}
\begin{proof}
  Write $a=2r$ and $b=2s$. Then $a+b=2(r+s)$.
\end{proof}
\begin{example}
  The integers $2$ and $4$ have even sum $6$.
\end{example}
Theorem~\ref{thm:even} gives the general statement.
\end{lstlisting}
\begin{definition}[Even integer]An integer $n$ is even if $n=2k$ for some $k\in\Z$.\end{definition}
\begin{theorem}\label{thm:even}The sum of two even integers is even.\end{theorem}
\begin{proof}Write $a=2r$ and $b=2s$. Then $a+b=2(r+s)$.\end{proof}
\begin{example}The integers $2$ and $4$ have even sum $6$.\end{example}
Theorem~\ref{thm:even} gives the general statement.
Theorem, definition, and example have separate counters and reset each
section. Lemma, proposition, corollary, and remark share the theorem counter.
A proof is unnumbered and adds its end-of-proof symbol automatically.

\section{Lists, tables, and images}
\subsection{Lists}
\begin{lstlisting}
\begin{enumerate}
  \item Check closure.
  \item Check inverses.
\end{enumerate}
\begin{itemize}
  \item A useful observation.
  \item Another observation.
\end{itemize}
\end{lstlisting}
\begin{enumerate}\item Check closure.\item Check inverses.\end{enumerate}
\begin{itemize}\item A useful observation.\item Another observation.\end{itemize}
Every item begins with \lstinline|\item|. Nest a list inside an item when you
need subitems; your \texttt{subproblem} environment gives Part A, Part B, etc.

\subsection{Tables}
\begin{lstlisting}
\begin{tabular}{lc}
  Object & Size \\
  \hline
  $\Z/2\Z$ & 2 \\
  $\Z/3\Z$ & 3
\end{tabular}
\end{lstlisting}
\begin{center}\begin{tabular}{lc}Object&Size\\\hline
$\Z/2\Z$&2\\$\Z/3\Z$&3\end{tabular}\end{center}
The column specifiers \texttt{l}, \texttt{c}, and \texttt{r} mean left,
center, and right. Ampersands separate columns; two backslashes end rows.
A \texttt{table} environment around \texttt{tabular} adds a floating caption
and label, like a figure.

\subsection{Images and captions}
Load \texttt{graphicx} in the preamble. Put an actual image at the referenced
path before compiling this example:
\begin{lstlisting}
\begin{figure}[htbp]
  \centering
  \includegraphics[width=0.65\linewidth]{images/diagram.png}
  \caption{A diagram of the construction.}
  \label{fig:construction}
\end{figure}
See Figure~\ref{fig:construction}.
\end{lstlisting}
A figure can move to a suitable page position; \texttt{htbp} gives placement
preferences. Put its label after the caption so references use the figure
number. This guide keeps the image code literal so it needs no external image.

\section{References and splitting large documents}
Choose descriptive, unique labels such as \texttt{eq:sum},
\texttt{thm:even}, and \texttt{sec:groups}. Use \lstinline|\ref{...}| for a
number and \lstinline|\eqref{...}| for an equation number with parentheses.
The tilde in \lstinline|Theorem~\ref{thm:even}| keeps the word and number
together. Compile again when references show question marks.
\begin{lstlisting}
% In the preamble:
\usepackage{hyperref}
% In the body:
\href{https://example.com}{Link text}
\url{https://example.com}
\end{lstlisting}

\subsection{Input files versus standalone lecture files}
For a simple split, keep a single preamble in the main file and insert body-only
files with \lstinline|\input{Lecture01}|. An input file should not repeat
\lstinline|\documentclass| or the document environment.
For lectures that can also compile alone, your \texttt{masterdoc} and
\texttt{lecturedoc} snippets use \texttt{subfiles} instead:
\begin{lstlisting}
% main.tex preamble:
\usepackage{subfiles}
% main.tex body:
\subfile{Lecture01}

% Lecture01.tex:
\documentclass[main.tex]{subfiles}
\begin{document}
\section{Lecture 1}
Lecture text.
\end{document}
\end{lstlisting}
The child inherits the main preamble. Its path to \texttt{main.tex} must be
correct. Your lecture template conditionally prints a standalone title so the
combined master does not get a second title page.

\section{Your homework workflow}
Load the style with the \texttt{solutions} option to show solutions. Set the
heading information before printing it:
\begin{lstlisting}
% Preamble:
\usepackage[solutions]{../../../mystyle/Style}
\sethwinfo{\defaultauthor}{Math 113}{Homework 8}
  {November 20, 2026}[None]
% Body:
\makehwtitle
\begin{problem}[3][10]
  Verify the group axioms.
\end{problem}
\begin{solution}
  \begin{subproblem}
    \item Check associativity.
    \item Find the identity.
    \item Find inverses.
  \end{subproblem}
\end{solution}
\end{lstlisting}
In the current style, the first optional problem argument overrides the printed
problem identifier; the second shows points. Without an override, numbering
advances automatically. The optional final argument of \lstinline|\sethwinfo|
sets collaborators and defaults to None. Without \texttt{solutions}, solution
content is hidden. For a quick check, compile once with the option and once
without it. Use a fixed due date instead of \lstinline|\today| when the date
should not change on recompilation.

\section{Your style commands: shortcut to expansion}
These are LaTeX macros: type the backslash command in source. They differ from
editor snippet triggers, which expand before compilation. Math commands in
this table belong inside math mode unless the note says otherwise.
\begin{longtable}{@{}p{0.27\linewidth}p{0.67\linewidth}@{}}
\toprule Command & Expansion or purpose \\\midrule\endhead
\lstinline|\N, \Z, \Q| & \lstinline|\mathbb{N}, \mathbb{Z}, \mathbb{Q}| \\
\lstinline|\R, \C, \F| & \lstinline|\mathbb{R}, \mathbb{C}, \mathbb{F}| \\
\lstinline|\abs{x}| & Absolute-value delimiters: \(\abs{x}\) \\
\lstinline|\norm{x}| & Norm delimiters: \(\norm{x}\) \\
\lstinline|\ceil{x}, \floor{x}| & Ceiling \(\ceil{x}\) and floor \(\floor{x}\) \\
\lstinline|\set{x,y}| & Braced set: \(\set{x,y}\) \\
\lstinline|\inner{x}{y}| & Inner product: \(\inner{x}{y}\) \\
\lstinline|\vect{v}| & \lstinline|\mathbf{v}|: \(\vect{v}\) \\
\lstinline|A\transpose| & \lstinline|A^{\mathsf{T}}|: \(A\transpose\) \\
\lstinline|\eps| & \lstinline|\varepsilon|: \(\eps\) \\
\lstinline|\phiinv| & \lstinline|\varphi^{-1}|: \(\phiinv\) \\
\lstinline|\defeq, \eqdef| & \lstinline|\vcentcolon=, =\vcentcolon|: \(x\defeq 3\), \(3\eqdef x\) \\
\lstinline|\given| & \lstinline|\mathrel{\middle\vert}|; use between matching left/right delimiters \\
\lstinline|\im, \rank, \Span| & Upright math operators: \(\im(f),\rank(A),\Span\set{u,v}\) \\
\lstinline|\Tr, \Var, \Cov| & Upright operators: \(\Tr(A),\Var(X),\Cov(X,Y)\) \\
\lstinline|\diag, \Res| & Upright operators: \(\diag(A),\Res(f)\) \\
\lstinline|\defaultauthor| & Your configured author name (text) \\
\lstinline|\sethwinfo| & Four required heading fields and optional collaborators (text) \\
\lstinline|\makehwtitle| & Print the homework heading (text) \\
\lstinline|\hwStudent, \hwCourse| & Current student and course heading fields (text) \\
\lstinline|\hwAssignment| & Current assignment field (text) \\
\lstinline|\hwDueDate| & Current due-date field (text) \\
\lstinline|\hwCollaborators| & Current collaborator field (text) \\
\lstinline|\TODO{Review}| & Red bold reminder (text) \\
\lstinline|\highlight{text}| & Yellow highlight; optional color, e.g. \lstinline|\highlight[green!25]{text}| \\
\lstinline|\marginref{Note}| & Small margin note (text) \\
\bottomrule
\end{longtable}
Paired-delimiter commands accept a star for automatic sizing, for example
\lstinline|\norm*{\frac{x}{y}}|. For conditional probability, use
\lstinline|\left( A \given B \right)| so \lstinline|\middle| has its
required matching delimiters. The raw heading fields can be changed with
\lstinline|\renewcommand|, but \lstinline|\sethwinfo| fills them together.

\section{Compiling and fixing errors}
Save the source and run your editor's LaTeX build or compile action. In VS Code,
use the LaTeX extension's build command if one is installed. A local terminal
alternative, run from the main file's folder, is:
\begin{lstlisting}
pdflatex -interaction=nonstopmode -halt-on-error main.tex
pdflatex -interaction=nonstopmode -halt-on-error main.tex
\end{lstlisting}
The second run resolves contents and references. In a multifile project,
compile the main file for the complete document.
\begin{description}
\item[Undefined control sequence.] Check spelling and whether the package or
style defining that command was loaded. For example, \lstinline|\R| needs
your style, and \lstinline|\includegraphics| needs \texttt{graphicx}.
\item[Missing dollar sign.] A math command or underscore may be outside math
mode, or a math delimiter may be unmatched. Inspect the preceding lines too.
\item[File not found.] Check the relative path, spelling, and capitalization
of a style, input file, or image. Do not add another random \texttt{../}.
\item[Environment ended by the wrong name.] Match begin/end names and braces.
\item[Misplaced alignment tab.] An ampersand belongs in a table or alignment;
write \lstinline|\&| for an ordinary ampersand.
\item[Question marks in references.] Compile again, then check that the label
exists, matches exactly, and belongs to a numbered object.
\item[Overfull box warning.] Content extends past a margin. Shorten an
unbreakable expression or split a long equation with an alignment.
\end{description}
Fix the first real error and compile again; later messages may be consequences
of that first error. Keep source and image files; auxiliary files such as
\texttt{.aux}, \texttt{.log}, and \texttt{.toc} are generated by compilation.

\clearpage
\section{All personal VS Code snippets}
\subsection{How to use the shortcuts}
In a LaTeX source file in VS Code, type a trigger such as \texttt{thm}, select
the matching snippet suggestion, and accept it. If suggestions are hidden,
use the editor's Insert Snippet command. Press Tab to move through editable
fields and Shift+Tab to move back. Tab completion behavior depends on your
editor settings; the trigger alone is not a LaTeX command.

The blocks below show the \emph{stored snippet bodies exactly}, after JSON
string decoding. Their placeholder notation is for the snippet engine:
\begin{itemize}
\item \lstinline|${1:Title}| means editable field 1, initially containing Title.
\item \lstinline|$1| and \lstinline|$2| are tab stops; repeated numbers share text.
\item \lstinline|$0| is the final cursor position and does not remain in TeX.
\end{itemize}
Thus \texttt{frac} becomes \lstinline|\frac{numerator}{denominator}| with
editable numerator and denominator fields. The \texttt{inline} body looks
unusual because it combines literal dollar signs with snippet placeholders:
after expansion it gives a single-dollar inline math pair, with the cursor
outside it. Do not paste raw placeholder bodies directly into a TeX document.
These are the 117 personal snippets from your supplied updated list;
extension-provided snippets and keyboard keybindings are separate.

\subsection{Document templates}
\subsubsection*{\texttt{latexdoc}: Basic LaTeX document}
\begin{lstlisting}
\documentclass{article}
\usepackage[utf8]{inputenc}

\title{${1:Title}}
\author{${2:Author}}
\date{\today}

\begin{document}
\maketitle

\section{${3:Introduction}}
$0

\end{document}
\end{lstlisting}
\subsubsection*{\texttt{mathdoc}: Math document using Style.sty}
\begin{lstlisting}
\documentclass{article}
\usepackage{../../../mystyle/Style}
\usepackage[margin=1in]{geometry}

\title{${1:Title}}
\author{\defaultauthor}
\date{\today}

\begin{document}
\maketitle

\section{${2:Introduction}}
$0

\end{document}
\end{lstlisting}
\subsubsection*{\texttt{basiclecture}: Standalone lecture document}
\begin{lstlisting}
\documentclass{article}
\usepackage{../../../mystyle/Style}
\usepackage[margin=1in]{geometry}

\title{$1 -- Lecture $2}
\author{\defaultauthor}
\date{$3}

\begin{document}
\maketitle

\section{$4}
$0

\end{document}
\end{lstlisting}
\subsubsection*{\texttt{masterdoc}: Master lecture document combining individual lectures}
\begin{lstlisting}
\documentclass[11pt]{article}
\usepackage{../../../mystyle/Style}
\usepackage[margin=1in]{geometry}
\usepackage[hidelinks]{hyperref}
\usepackage{subfiles}

\title{${1:Course}: ${2:Title}}
\author{\defaultauthor}
\date{${3:Semester}}

\begin{document}
\maketitle
\tableofcontents
\newpage

% Add lectures below
\subfile{${4:Lecture1}}
$0

\end{document}
\end{lstlisting}
\subsubsection*{\texttt{lecturedoc}: Lecture with standalone title and master document support}
\begin{lstlisting}
\documentclass[main.tex]{subfiles}

% Individual lecture title
\ifSubfilesClassLoaded{%
    \title{${1:Course}: ${2:Topic}}
    \author{\defaultauthor}
    \date{${3:Lecture Date}}
}{}

\begin{document}

\ifSubfilesClassLoaded{\maketitle}{}

\section{Lecture ${4:Number}: ${2:Topic}}

$0

\end{document}
\end{lstlisting}
\subsubsection*{\texttt{hwdoc}: Homework using Style.sty and a fixed due date}
\begin{lstlisting}
\documentclass{article}
\usepackage[solutions]{../../../mystyle/Style}
\usepackage[margin=1in]{geometry}

\sethwinfo{\defaultauthor}{${1:Course}}{${2:Homework}}{${3:Due Date}}

\begin{document}
\makehwtitle

\begin{problem}
	${4:Problem statement}
\end{problem}

\begin{solution}
	$0
\end{solution}

\end{document}
\end{lstlisting}

\subsection{Theorem and homework environments}
\subsubsection*{\texttt{thm}: Theorem environment}
\begin{lstlisting}
\begin{theorem}
	$0
\end{theorem}
\end{lstlisting}
\subsubsection*{\texttt{def}: Definition environment}
\begin{lstlisting}
\begin{definition}
	$0
\end{definition}
\end{lstlisting}
\subsubsection*{\texttt{lem}: Lemma environment}
\begin{lstlisting}
\begin{lemma}
	$0
\end{lemma}
\end{lstlisting}
\subsubsection*{\texttt{prop}: Proposition environment}
\begin{lstlisting}
\begin{proposition}
	$0
\end{proposition}
\end{lstlisting}
\subsubsection*{\texttt{cor}: Corollary environment}
\begin{lstlisting}
\begin{corollary}
	$0
\end{corollary}
\end{lstlisting}
\subsubsection*{\texttt{ex}: Example environment}
\begin{lstlisting}
\begin{example}
	$0
\end{example}
\end{lstlisting}
\subsubsection*{\texttt{rem}: Remark environment}
\begin{lstlisting}
\begin{remark}
	$0
\end{remark}
\end{lstlisting}
\subsubsection*{\texttt{prf}: Proof environment}
\begin{lstlisting}
\begin{proof}
	$0
\end{proof}
\end{lstlisting}
\subsubsection*{\texttt{prob}: Homework problem environment}
\begin{lstlisting}
\begin{problem}
	$0
\end{problem}
\end{lstlisting}
\subsubsection*{\texttt{subprob}: Homework subproblem environment}
\begin{lstlisting}
\begin{subproblem}
	\item $0
\end{subproblem}
\end{lstlisting}
\subsubsection*{\texttt{sol}: Homework solution environment}
\begin{lstlisting}
\begin{solution}
	$0
\end{solution}
\end{lstlisting}

\subsection{Equations and basic math}
\subsubsection*{\texttt{inline}: Inline math}
\begin{lstlisting}
$${1:expression}$$0
\end{lstlisting}
\subsubsection*{\texttt{display}: Display math}
\begin{lstlisting}
\[
	$0
\]
\end{lstlisting}
\subsubsection*{\texttt{eq}: Unnumbered equation}
\begin{lstlisting}
\begin{equation*}
	$0
\end{equation*}
\end{lstlisting}
\subsubsection*{\texttt{align}: Unnumbered aligned equations}
\begin{lstlisting}
\begin{align*}
	$0
\end{align*}
\end{lstlisting}
\subsubsection*{\texttt{frac}: Fraction}
\begin{lstlisting}
\frac{${1:numerator}}{${2:denominator}}$0
\end{lstlisting}
\subsubsection*{\texttt{sqrt}: Square root}
\begin{lstlisting}
\sqrt{${1:expression}}$0
\end{lstlisting}

\subsection{Calculus and sets}
\subsubsection*{\texttt{sum}: Summation}
\begin{lstlisting}
\sum_{${1:i=1}}^{${2:n}} $0
\end{lstlisting}
\subsubsection*{\texttt{prod}: Product}
\begin{lstlisting}
\prod_{${1:i=1}}^{${2:n}} $0
\end{lstlisting}
\subsubsection*{\texttt{lim}: Limit}
\begin{lstlisting}
\lim_{${1:n}\to ${2:\infty}} $0
\end{lstlisting}
\subsubsection*{\texttt{int}: Definite integral}
\begin{lstlisting}
\int_{${1:a}}^{${2:b}} ${3:f(x)}\,d${4:x}$0
\end{lstlisting}
\subsubsection*{\texttt{set}: Set using Style.sty}
\begin{lstlisting}
\set{${1:elements}}$0
\end{lstlisting}
\subsubsection*{\texttt{setb}: Set-builder notation}
\begin{lstlisting}
\set{${1:x} \mid ${2:condition}}$0
\end{lstlisting}
\subsubsection*{\texttt{abs}: Absolute value using Style.sty}
\begin{lstlisting}
\abs{${1:x}}$0
\end{lstlisting}
\subsubsection*{\texttt{norm}: Norm using Style.sty}
\begin{lstlisting}
\norm{${1:x}}$0
\end{lstlisting}
\subsubsection*{\texttt{ceil}: Ceiling using Style.sty}
\begin{lstlisting}
\ceil{${1:x}}$0
\end{lstlisting}
\subsubsection*{\texttt{floor}: Floor using Style.sty}
\begin{lstlisting}
\floor{${1:x}}$0
\end{lstlisting}
\subsubsection*{\texttt{inner}: Inner product using Style.sty}
\begin{lstlisting}
\inner{${1:x}}{${2:y}}$0
\end{lstlisting}

\subsection{Text and math styling}
These shortcuts wrap selected text. With no selection, type into the empty
field; for \texttt{col;}, choose the color first (red by default), then the text.
The \lstinline|${TM_SELECTED_TEXT:$1}| variable uses the selection when
available and otherwise provides an editable tab stop. Press Tab to reach
\lstinline|$0| after the closing brace.
Use \texttt{mbf;} and \texttt{bsy;} inside math mode: \lstinline|\mathbf|
bolds mathematical letters, while \lstinline|\boldsymbol| also bolds symbols
such as Greek letters. Colored text requires \texttt{xcolor}, and
\lstinline|\boldsymbol| requires \texttt{amsmath}; your style file loads both.

\subsubsection*{\texttt{bf;}: Bold text}
\begin{lstlisting}
\textbf{${TM_SELECTED_TEXT:$1}}$0
\end{lstlisting}
\subsubsection*{\texttt{it;}: Italic text}
\begin{lstlisting}
\textit{${TM_SELECTED_TEXT:$1}}$0
\end{lstlisting}
\subsubsection*{\texttt{em;}: Emphasized text}
\begin{lstlisting}
\emph{${TM_SELECTED_TEXT:$1}}$0
\end{lstlisting}
\subsubsection*{\texttt{ul;}: Underlined text}
\begin{lstlisting}
\underline{${TM_SELECTED_TEXT:$1}}$0
\end{lstlisting}
\subsubsection*{\texttt{col;}: Colored text}
\begin{lstlisting}
\textcolor{${1:red}}{${TM_SELECTED_TEXT:$2}}$0
\end{lstlisting}
\subsubsection*{\texttt{mbf;}: Bold mathematical letters}
\begin{lstlisting}
\mathbf{${TM_SELECTED_TEXT:$1}}$0
\end{lstlisting}
\subsubsection*{\texttt{bsy;}: Bold mathematical symbols}
\begin{lstlisting}
\boldsymbol{${TM_SELECTED_TEXT:$1}}$0
\end{lstlisting}

\subsection{Matrices and math formatting}
\subsubsection*{\texttt{mat}: 2x2 parenthesis matrix}
\begin{lstlisting}
\begin{pmatrix}
	${1:a} & ${2:b} \\
	${3:c} & ${4:d}
\end{pmatrix}$0
\end{lstlisting}
\subsubsection*{\texttt{bmat}: 2x2 bracket matrix}
\begin{lstlisting}
\begin{bmatrix}
	${1:a} & ${2:b} \\
	${3:c} & ${4:d}
\end{bmatrix}$0
\end{lstlisting}
\subsubsection*{\texttt{cases}: Piecewise cases}
\begin{lstlisting}
\begin{cases}
	${1:expression} & ${2:condition} \\
	${3:expression} & ${4:condition}
\end{cases}$0
\end{lstlisting}
\subsubsection*{\texttt{paren}: Automatically sized parentheses}
\begin{lstlisting}
\left(${1:expression}\right)$0
\end{lstlisting}
\subsubsection*{\texttt{bracket}: Automatically sized brackets}
\begin{lstlisting}
\left[${1:expression}\right]$0
\end{lstlisting}
\subsubsection*{\texttt{text}: Text inside math mode}
\begin{lstlisting}
\text{${1:text}}$0
\end{lstlisting}
\subsubsection*{\texttt{bb}: Blackboard bold}
\begin{lstlisting}
\mathbb{${1:R}}$0
\end{lstlisting}
\subsubsection*{\texttt{cal}: Calligraphic letters}
\begin{lstlisting}
\mathcal{${1:F}}$0
\end{lstlisting}
\subsubsection*{\texttt{frak}: Fraktur letters}
\begin{lstlisting}
\mathfrak{${1:g}}$0
\end{lstlisting}

\subsection{Document organization}
\subsubsection*{\texttt{env}: Generic LaTeX environment}
\begin{lstlisting}
\begin{${1:environment}}
	$0
\end{${1:environment}}
\end{lstlisting}
\subsubsection*{\texttt{enum}: Numbered list}
\begin{lstlisting}
\begin{enumerate}
	\item $0
\end{enumerate}
\end{lstlisting}
\subsubsection*{\texttt{item}: Bullet list}
\begin{lstlisting}
\begin{itemize}
	\item $0
\end{itemize}
\end{lstlisting}
\subsubsection*{\texttt{sec}: Section heading}
\begin{lstlisting}
\section{${1:Title}}$0
\end{lstlisting}
\subsubsection*{\texttt{subsec}: Subsection heading}
\begin{lstlisting}
\subsection{${1:Title}}$0
\end{lstlisting}
\subsubsection*{\texttt{input}: Insert contents of another TeX file}
\begin{lstlisting}
\input{${1:filename}}$0
\end{lstlisting}


\subsection{Symbol shortcuts: trigger, command, and result}
These snippets insert math symbols, so expand them inside inline or display
math. Triggers are case-sensitive: \lstinline!;g! inserts
\lstinline!\gamma!, while \lstinline!;G! inserts \lstinline!\Gamma!.
The semicolon belongs to the editor trigger and is replaced during expansion.
For example, expanding \lstinline!;a! inside \lstinline!$...$! inserts
\lstinline!\alpha! and renders as \(\alpha\).

Your default epsilon and phi triggers use \lstinline!\varepsilon! and
\lstinline!\varphi!; the alternate triggers give \lstinline!\epsilon! and
\lstinline!\phi!. Omicron inserts the ordinary letter \lstinline!o!.
The \lstinline!iso! shortcut inserts \lstinline!\cong! (\(\cong\)).

\subsubsection*{Lowercase Greek letters}
\begin{longtable}{@{}p{0.12\linewidth}p{0.25\linewidth}p{0.11\linewidth}p{0.42\linewidth}@{}}
\toprule Trigger & Inserts & Result & Meaning \\ \midrule
\endhead
\lstinline!;a! & \lstinline!\alpha! & \(\alpha\) & Greek alpha \\
\lstinline!;b! & \lstinline!\beta! & \(\beta\) & Greek beta \\
\lstinline!;g! & \lstinline!\gamma! & \(\gamma\) & Greek gamma \\
\lstinline!;d! & \lstinline!\delta! & \(\delta\) & Greek delta \\
\lstinline!;e! & \lstinline!\varepsilon! & \(\varepsilon\) & Greek epsilon \\
\lstinline!;z! & \lstinline!\zeta! & \(\zeta\) & Greek zeta \\
\lstinline!;h! & \lstinline!\eta! & \(\eta\) & Greek eta \\
\lstinline!;t! & \lstinline!\theta! & \(\theta\) & Greek theta \\
\lstinline!;i! & \lstinline!\iota! & \(\iota\) & Greek iota \\
\lstinline!;k! & \lstinline!\kappa! & \(\kappa\) & Greek kappa \\
\lstinline!;l! & \lstinline!\lambda! & \(\lambda\) & Greek lambda \\
\lstinline!;m! & \lstinline!\mu! & \(\mu\) & Greek mu \\
\lstinline!;n! & \lstinline!\nu! & \(\nu\) & Greek nu \\
\lstinline!;x! & \lstinline!\xi! & \(\xi\) & Greek xi \\
\lstinline!;o! & \lstinline!o! & \(o\) & Greek omicron \\
\lstinline!;p! & \lstinline!\pi! & \(\pi\) & Greek pi \\
\lstinline!;r! & \lstinline!\rho! & \(\rho\) & Greek rho \\
\lstinline!;s! & \lstinline!\sigma! & \(\sigma\) & Greek sigma \\
\lstinline!;u! & \lstinline!\tau! & \(\tau\) & Greek tau \\
\lstinline!;y! & \lstinline!\upsilon! & \(\upsilon\) & Greek upsilon \\
\lstinline!;f! & \lstinline!\varphi! & \(\varphi\) & Greek phi \\
\lstinline!;c! & \lstinline!\chi! & \(\chi\) & Greek chi \\
\lstinline!;j! & \lstinline!\psi! & \(\psi\) & Greek psi \\
\lstinline!;w! & \lstinline!\omega! & \(\omega\) & Greek omega \\
\bottomrule
\end{longtable}

\subsubsection*{Alternate Greek letters}
\begin{longtable}{@{}p{0.12\linewidth}p{0.25\linewidth}p{0.11\linewidth}p{0.42\linewidth}@{}}
\toprule Trigger & Inserts & Result & Meaning \\ \midrule
\endhead
\lstinline!;ve! & \lstinline!\epsilon! & \(\epsilon\) & Greek alternate epsilon \\
\lstinline!;vf! & \lstinline!\phi! & \(\phi\) & Greek alternate phi \\
\lstinline!;vt! & \lstinline!\vartheta! & \(\vartheta\) & Greek alternate theta \\
\lstinline!;vr! & \lstinline!\varrho! & \(\varrho\) & Greek alternate rho \\
\lstinline!;vs! & \lstinline!\varsigma! & \(\varsigma\) & Greek alternate sigma \\
\lstinline!;vp! & \lstinline!\varpi! & \(\varpi\) & Greek alternate pi \\
\lstinline!;vk! & \lstinline!\varkappa! & \(\varkappa\) & Greek alternate kappa \\
\bottomrule
\end{longtable}

\subsubsection*{Uppercase Greek letters}
\begin{longtable}{@{}p{0.12\linewidth}p{0.25\linewidth}p{0.11\linewidth}p{0.42\linewidth}@{}}
\toprule Trigger & Inserts & Result & Meaning \\ \midrule
\endhead
\lstinline!;G! & \lstinline!\Gamma! & \(\Gamma\) & Greek Gamma \\
\lstinline!;D! & \lstinline!\Delta! & \(\Delta\) & Greek Delta \\
\lstinline!;T! & \lstinline!\Theta! & \(\Theta\) & Greek Theta \\
\lstinline!;L! & \lstinline!\Lambda! & \(\Lambda\) & Greek Lambda \\
\lstinline!;X! & \lstinline!\Xi! & \(\Xi\) & Greek Xi \\
\lstinline!;P! & \lstinline!\Pi! & \(\Pi\) & Greek Pi \\
\lstinline!;S! & \lstinline!\Sigma! & \(\Sigma\) & Greek Sigma \\
\lstinline!;Y! & \lstinline!\Upsilon! & \(\Upsilon\) & Greek Upsilon \\
\lstinline!;F! & \lstinline!\Phi! & \(\Phi\) & Greek Phi \\
\lstinline!;J! & \lstinline!\Psi! & \(\Psi\) & Greek Psi \\
\lstinline!;W! & \lstinline!\Omega! & \(\Omega\) & Greek Omega \\
\bottomrule
\end{longtable}

\subsubsection*{Arrows}
\begin{longtable}{@{}p{0.12\linewidth}p{0.25\linewidth}p{0.11\linewidth}p{0.42\linewidth}@{}}
\toprule Trigger & Inserts & Result & Meaning \\ \midrule
\endhead
\lstinline!->! & \lstinline!\to! & \(\to\) & Right Arrow \\
\lstinline!<-! & \lstinline!\leftarrow! & \(\leftarrow\) & Left Arrow \\
\lstinline!<->! & \lstinline!\leftrightarrow! & \(\leftrightarrow\) & Bidirectional Arrow \\
\lstinline!=>! & \lstinline!\implies! & \(\implies\) & Implies \\
\lstinline!<=>! & \lstinline!\iff! & \(\iff\) & If and Only If \\
\lstinline!|->! & \lstinline!\mapsto! & \(\mapsto\) & Maps To \\
\lstinline!up! & \lstinline!\uparrow! & \(\uparrow\) & Up Arrow \\
\lstinline!down! & \lstinline!\downarrow! & \(\downarrow\) & Down Arrow \\
\bottomrule
\end{longtable}

\subsubsection*{Mathematical symbols}
\begin{longtable}{@{}p{0.12\linewidth}p{0.25\linewidth}p{0.11\linewidth}p{0.42\linewidth}@{}}
\toprule Trigger & Inserts & Result & Meaning \\ \midrule
\endhead
\lstinline!inn! & \lstinline!\in! & \(\in\) & Element Of \\
\lstinline!nin! & \lstinline!\notin! & \(\notin\) & Not Element Of \\
\lstinline!sub! & \lstinline!\subseteq! & \(\subseteq\) & Subset Or Equal \\
\lstinline!sup! & \lstinline!\supseteq! & \(\supseteq\) & Superset Or Equal \\
\lstinline!forall! & \lstinline!\forall! & \(\forall\) & For All \\
\lstinline!exists! & \lstinline!\exists! & \(\exists\) & There Exists \\
\lstinline!neq! & \lstinline!\neq! & \(\neq\) & Not Equal \\
\lstinline!leq! & \lstinline!\leq! & \(\leq\) & Less Than Or Equal \\
\lstinline!geq! & \lstinline!\geq! & \(\geq\) & Greater Than Or Equal \\
\lstinline!iso! & \lstinline!\cong! & \(\cong\) & Isomorphic \\
\lstinline!empty! & \lstinline!\varnothing! & \(\varnothing\) & Empty Set \\
\bottomrule
\end{longtable}

\end{document}
